\documentclass[10pt,a4paper]{amsart}
\usepackage[margin=19mm]{geometry}
\usepackage{amsmath,amssymb,mathtools}
\usepackage[scr=rsfs,cal=euler]{mathalfa}
\usepackage{xfrac}
\usepackage[unicode,hidelinks,bookmarks=false]{hyperref}
\newcommand{\ie}{i.e.\ }
\newcommand{\cf}{cf.\ }
\newcommand{\resp}{respectively\ }
\newcommand{\Subsection}{\subsection}
\providecommand{\nobreakdash}{\nobreak\hbox{-}}
\setlength{\emergencystretch}{2em}
\multlinegap0pt
\usepackage{xcolor}
\usepackage{amssymb}
\newtheorem{lemma}{Lemma}
\title{Optimal Union Probability Interval Is NP-Hard}
\author{Petteri Kaski, Heikki Mannila, Chandra Kanta Mohapatra}
\date{Source: arXiv:2605.03556v2 (CC BY 4.0)}
\begin{document}
\maketitle

\noindent\emph{Source and licence.} Optimal Union Probability Interval Is NP-Hard. Version arXiv:2605.03556v2 (CC BY 4.0), distributed by arXiv under the Creative Commons Attribution 4.0 International licence, http://creativecommons.org/licenses/by/4.0/, as recorded in the arXiv licence metadata for this exact version and shown on the abstract page https://arxiv.org/abs/2605.03556v2. Full licence notice: \texttt{licenses/arxiv-2605.03556-CC-BY-4.0.txt}.\par\smallskip

\noindent\textit{Editorial scope.} Counted unit: the article's lemma lem:special-dual (source0.tex lines 730--749). The article's complete original proof of it is delivered with it (source0.tex lines 751--897, one \texttt{proof} environment of 147 lines). No mathematics is written, repaired or re-derived: the statement and the proof are the article's own lines, in the article's own notation, and no retained line is substituted or re-flowed.  No further source-local context is needed: every cross-reference of every retained line resolves inside the two delivered spans.  Nothing was omitted from any retained span, so the package carries no omission record.  No retained line contains a citation, so the wrapper carries no bibliography and no bibliography entry.  No retained line contains a figure environment, an image-inclusion command or drawing code, so no image is delivered and the package declares no asset dependency.  Editorial formatting only, in addition: an AMS \texttt{amsart} preamble that restates the article's own declarations and macros used by the retained lines, namely the environments \texttt{lemma} on separate counters, together with the standard packages those lines need.  The retained environments print in this document as Lemma 1 in order of appearance, whereas the article prints the same items with the numbering of its own full document.  The complete native source of the article as distributed in the arXiv e-print for this exact version is bundled unchanged as \texttt{sources/199797/source0.tex}.
\begin{lemma}[A transformed special case of the dual linear program]
\label{lem:special-dual}
Assume access to an oracle that solves the maximum union
probability problem for a set family $\emptyset\neq\mathcal{F}\subseteq2^{[n]}$ 
and a feasible rational $b:\mathcal{F}\rightarrow [0,1]$ given as query.
Then, given a nonempty $n$-vertex graph $G$ and edge weights
$w:E(G)\rightarrow\mathbb{Q}_{\geq 0}$
as input, we can solve the linear program to
\begin{equation}
\label{eq:dual-simplified}
\begin{split}
&\text{maximize $\sum_{\{u,v\}\in E(G)}w_{\{u,v\}}y_{\{u,v\}}$}\\
&\text{subject to $\sum_{E(G)\ni \{u,v\}\subseteq T}y_{\{u,v\}}\leq |T|-1$ 
for all cliques $\emptyset\neq T\subseteq V(G)$}
\end{split}
\end{equation}
for its optimum value in time polynomial in $n$ and the rational encoding 
length of $w$, by making a single query to the maximum union probability 
oracle, with $\mathcal{F}$ consisting of sets of size at most two.
\end{lemma}

\begin{proof}
Let us recall the maximum union probability problem as a linear program 
on a given query $(b,\mathcal{F})$ with $b$ feasible. Namely, we are to
\begin{equation}
\label{eq:max-union-lp}
\begin{split}
&\text{maximize $\sum_{\emptyset\neq T\subseteq[n]}x_T$}\\
&\text{subject to $\sum_{T\subseteq[n]}x_T=1$, $\sum_{S\subseteq T}x_T=b_S$ for all $S\in\mathcal{F}$, and $x_T\geq 0$ for all $T\subseteq[n]$}\,. 
\end{split}
\end{equation}
We assume we have an oracle subroutine that on query
$(b,\mathcal{F})$ with $b$ feasible 
outputs the value of \eqref{eq:max-union-lp}.

Let $G$ be a graph and $w:E(G)\rightarrow\mathbb{Q}_{\geq 0}$ a function
given as input. (We will discuss the input $w$ in detail only towards
the end of the proof of this lemma.) 
We may assume that $n\geq 2$ and that $G$ has at least one edge.
We may also assume that the vertex set $V(G)=[n]$ and that the edge
set $E(G)$ consists of $2$-element subsets of $V(G)$.
Fix
\[
\mathcal{F}=\{\{u\}:u\in V(G)\}\cup\{\{u,v\}:u,v\in V(G)\}\,.
\]
We will also fix $b_{\{u\}}=\frac{1}{n}$ and $b_{\{u,v\}}=0$ 
for all distinct $u,v\in V(G)$ with $\{u,v\}\notin E(G)$ in our
query to~\eqref{eq:max-union-lp} in what follows. 
With this fixing, the query $(b,\mathcal{F})$ 
to~\eqref{eq:max-union-lp} that we will make in what follows 
is defined by a function $c:E(G)\rightarrow[0,\frac{1}{n^2}]$ when we set
$b_{\{u,v\}}=c_{\{u,v\}}$ for all $\{u,v\}\in E(G)$; the upper bound 
$\frac{1}{n^2}$ for the values of $c$ is to ensure feasibility of the
resulting query $b$. Indeed, to see that such a function $c$ results 
in a feasible query $b$ to~\eqref{eq:max-union-lp}, observe that $b$ 
is realized by the atom probabilities $x$ defined for all $T\subseteq[n]$ by
\[
x_T=\begin{cases}
0 & \text{if $|T|>2$, or $|T|=2$ and $T\notin E(G)$},\\
c_T & \text{if $|T|=2$ and $T\in E(G)$},\\
\frac{1}{n}-\sum_{T\subsetneq U\subseteq[n]}x_U & \text{if $|T|=1$},\\
1-\sum_{\emptyset\neq U\subseteq[n]}x_U & \text{if $T=\emptyset$}.
\end{cases}
\]
Now let us assume a function $c:E(G)\rightarrow[0,\frac{1}{n^2}]$ and
thus a query $(b,\mathcal{F})$ has been fixed. We will proceed along a
sequence of linear programs equivalent to \eqref{eq:max-union-lp}. 
Since we have 
fixed $b_{\{u\}}=\frac{1}{n}$ for all $u\in V(G)$ and $x_T\geq 0$ holds 
for all $\emptyset\neq T\subseteq[n]$, we observe that 
\[
\sum_{\emptyset\neq T\subseteq[n]}x_T\leq \sum_{u\in V(G)}\sum_{u\in T\subseteq V(G)}x_T\leq n\cdot\frac{1}{n}=1\,. 
\]
This inequality in particular implies that we can remove the constraint
$\sum_{T\subseteq[n]}x_T=1$ and the variable $x_\emptyset$ 
from \eqref{eq:max-union-lp}; that is, an equivalent linear program is to
\begin{equation}
\label{eq:max-union-lp-trim}
\begin{split}
&\text{maximize $\sum_{\emptyset\neq T\subseteq[n]}x_T$}\\
&\text{subject to $\sum_{S\subseteq T}x_T=b_S$ for all $S\in\mathcal{F}$
and $x_T\geq 0$ for all $\emptyset\neq T\subseteq[n]$}.
\end{split}
\end{equation}
The dual of \eqref{eq:max-union-lp-trim} is to
\begin{equation}
\label{eq:max-union-lp-trim-dual}
\begin{split}
&\text{minimize $\sum_{S\in\mathcal{F}}b_Sy_S$}\\
&\text{subject to $\sum_{\mathcal{F}\ni S\subseteq T}y_S\geq 1$ for all $\emptyset\neq T\subseteq[n]$},
\end{split}
\end{equation}
where each variable $y_S$ for $S\in\mathcal{F}$ can take arbitrary rational
values. Since $b_{\{u,v\}}=0$ whenever $\{u,v\}\notin E(G)$, we observe
that in \eqref{eq:max-union-lp-trim-dual} each constraint for 
$\emptyset\neq T\subseteq[n]$ that is not a clique in $G$ can be 
trivially satisfied by assigning an arbitrarily large value to a 
variable $y_{\{u,v\}}$ for $T\supseteq\{u,v\}\notin E(G)$, 
without affecting the value of the objective function since $b_{\{u,v\}}=0$.
In particular, we can remove all the non-clique constraints and all the
variables $y_{\{u,v\}}$ with $\{u,v\}\notin E(G)$ 
from \eqref{eq:max-union-lp-trim-dual} to obtain the equivalent linear 
program to 
\begin{equation}
\label{eq:max-union-lp-trim-dual-clique}
\begin{split}
&\text{minimize $\frac{1}{n}\sum_{u\in V(G)}y_{\{u\}}+\sum_{\{u,v\}\in E(G)}c_{\{u,v\}}y_{\{u,v\}}$}\\
&\text{subject to $\sum_{V(G)\ni u\in T}y_{\{u\}}+\sum_{E(G)\ni \{u,v\}\subseteq T}y_{\{u,v\}}\geq 1$ for all cliques $\emptyset\neq T\subseteq V(G)$}.
\end{split}
\end{equation}
Now fix an arbitrary optimum solution $y^*$ to \eqref{eq:max-union-lp-trim-dual-clique}. Without loss of generality we can assume that 
$y^*_{\{u\}}=1$ holds for all $u\in V(G)$; 
indeed, $y^*_{\{u\}}\geq 1$ is immediate by the clique $T=\{u\}$, and 
the assumption $y^*_{\{u_0\}}=1+\epsilon$ for some $u_0\in V(G)$
and $\epsilon>0$ yields a contradiction to the optimality of 
$y^*$ by decreasing $y^*_{\{u_0\}}$ (with objective coefficient $\frac{1}{n}$) 
by $\epsilon$ and increasing each $y^*_{\{u_0,v\}}$ 
(with objective coefficient $c_{\{u_0,v\}}\leq\frac{1}{n^2}$)
by $\epsilon$ for each edge $\{u_0,v\}\in E(G)$ 
to obtain an overall decrease in the minimization objective of 
\eqref{eq:max-union-lp-trim-dual-clique} without 
violating any of the clique constraints.
Thus, a linear program equivalent to \eqref{eq:max-union-lp-trim-dual-clique}
is to
\begin{equation}
\label{eq:dual-simplified-1}
\begin{split}
&\text{minimize $1+\sum_{\{u,v\}\in E(G)}c_{\{u,v\}}y_{\{u,v\}}$}\\
&\text{subject to $|T|+\sum_{E(G)\ni \{u,v\}\subseteq T}y_{\{u,v\}}\geq 1$ 
for all cliques $\emptyset\neq T\subseteq V(G)$}. 
\end{split}
\end{equation}
Recalling that the variables 
$y_{\{u,v\}}$ can take arbitrary rational values, changing the sign 
of the variables and removing the known offset $1$ from the objective, 
we observe that an offset-$1$-equivalent linear program 
to \eqref{eq:dual-simplified-1} is to 
\begin{equation}
\label{eq:dual-simplified-2}
\begin{split}
&\text{maximize $\sum_{\{u,v\}\in E(G)}c_{\{u,v\}}y_{\{u,v\}}$}\\
&\text{subject to $\sum_{E(G)\ni \{u,v\}\subseteq T}y_{\{u,v\}}\leq |T|-1$ 
for all cliques $\emptyset\neq T\subseteq V(G)$}.
\end{split}
\end{equation}
It follows immediately from this sequence \eqref{eq:max-union-lp}--\eqref{eq:dual-simplified-2} of offset-equivalent linear programs
that we can use our assumed maximum union probability oracle to obtain 
the maximum value of~\eqref{eq:dual-simplified} 
for an arbitrary $c:E(G)\rightarrow [0,\frac{1}{n^2}]$.

Now consider the given rational input $w:E(G)\rightarrow \mathbb{Q}_{\geq 0}$
and recall that we are to solve~\eqref{eq:dual-simplified} for its
optimum value. 
When $w$ is identically zero, the optimum value of~\eqref{eq:dual-simplified} 
is zero, so we can assume that $\max_{e\in E(G)}w(e)>0$ in what follows. 
Set $\alpha=(n^2\max_{e\in E(G)}w_e)^{-1}$ and $c=\alpha w$, observing that
all values of $c$ are rational in the interval $[0,\frac{1}{n^2}]$. 
Furthermore, $c$ has rational encoding length at most 
polynomial in $n$ and the rational encoding length of~$w$. 
Now invoke the maximum union probability oracle on input $(b,\mathcal{F})$
derived from $c$ and offset the oracle output by~$1$ to obtain the optimum 
value $t$ for~\eqref{eq:dual-simplified-2}, and return the value 
$t/\alpha$, observing that the returned value is the optimum value
of~\eqref{eq:dual-simplified} on input $w$. This process runs in time
bounded by a polynomial in $n$ and the rational encoding length of $w$,
making a single query $(b,\mathcal{F})$ to the maximum union probability 
oracle with $\mathcal{F}$ consisting of sets of size at most two. 
\end{proof}

\end{document}
